\chapter*{Conclusion Générale}
\addcontentsline{toc}{chapter}{Conclusion Générale}

\noindent
La réalisation de \projectName{} a permis de mettre en pratique les compétences acquises en ingénierie logicielle, notamment dans la conception d'une application web complète, l'intégration d'une API REST, la gestion d'une base de données relationnelle et la mise en place de parcours utilisateurs différenciés selon les rôles.

\noindent
Le projet répond à un besoin concret du secteur CHR : mieux organiser le recrutement, centraliser les offres et les candidatures, réutiliser le CV du candidat, qualifier les profils par quiz et offrir aux recruteurs un espace de suivi professionnel. L'ajout d'un espace administrateur permet également de contrôler les offres publiées et de rendre la plateforme plus fiable.

\noindent
Sur le plan technique, ce stage m'a permis de renforcer mes compétences en React, Laravel, MySQL, authentification par token, validation des fichiers, gestion des rôles, tests backend et amélioration de l'expérience utilisateur. Il m'a aussi permis de comprendre l'importance de la stabilisation finale d'un projet avant une démonstration : cohérence des boutons, navigation, messages d'erreur, états vides et règles métier.

\noindent
Comme perspectives, plusieurs améliorations peuvent être envisagées : intégration complète d'un paiement réel, notifications par email, statistiques plus avancées, amélioration de l'algorithme de matching, application mobile et intégration éventuelle avec d'autres solutions de gestion existantes. Ces évolutions permettraient de transformer \projectName{} en une plateforme encore plus complète et adaptée aux besoins réels des établissements CHR.
